\chapter{Objetivos y metodología}
\label{cap:objetivos}

\section{Objetivo general}
El objetivo general de este Trabajo de Fin de Grado es \textbf{investigar la
aplicación del Análisis Topológico de Datos (TDA) a la clasificación y
visualización de mamografías digitales}, evaluando su capacidad discriminativa
entre lesiones benignas y malignas y su complementariedad con los modelos de
aprendizaje profundo.

\section{Objetivos específicos}
Para alcanzar el objetivo general se plantean los siguientes objetivos
específicos:

\begin{enumerate}
  \item[\textbf{O1.}] \textbf{Revisar el estado del arte} en (i) análisis
        automático de mamografías y (ii) aplicaciones del TDA en imagen médica.
  \item[\textbf{O2.}] \textbf{Preparar y preprocesar} el conjunto de datos
        \mbox{CBIS-DDSM}: descarga, limpieza de metadatos, extracción de ROIs y
        normalización de las imágenes.
  \item[\textbf{O3.}] \textbf{Extraer descriptores topológicos} (diagramas de
        persistencia, códigos de barras e imágenes de persistencia) a partir de
        las mamografías mediante homología persistente.
  \item[\textbf{O4.}] \textbf{Entrenar y evaluar clasificadores} sobre los
        descriptores topológicos (p.\,ej. SVM, \textit{random forest}) para la
        tarea benigno/maligno.
  \item[\textbf{O5.}] \textbf{Comparar} el enfoque topológico con un modelo de
        \textit{deep learning} de referencia (CNN con \textit{transfer
        learning}) y \textbf{explorar su combinación} (fusión de
        características).
  \item[\textbf{O6.}] \textbf{Desarrollar herramientas de visualización} que
        hagan interpretables los resultados topológicos para su uso como apoyo
        al diagnóstico.
\end{enumerate}

\section{Metodología general}
El desarrollo del trabajo sigue una metodología incremental e iterativa,
estructurada en las siguientes fases, alineadas con los objetivos anteriores:

\begin{enumerate}
  \item \textbf{Fase de documentación} (O1): estudio de la bibliografía y de las
        herramientas disponibles.
  \item \textbf{Fase de datos} (O2): construcción de un \textit{pipeline}
        reproducible de preparación de datos.
  \item \textbf{Fase de extracción de características} (O3): implementación de la
        extracción de descriptores topológicos.
  \item \textbf{Fase de modelado} (O4, O5): entrenamiento, validación y
        comparación de modelos, evitando la fuga de datos (\textit{data
        leakage}) mediante una separación estricta \textit{train/validation/test}.
  \item \textbf{Fase de análisis y visualización} (O6): interpretación de los
        resultados y desarrollo de visualizaciones.
  \item \textbf{Fase de redacción} de la memoria.
\end{enumerate}

\begin{table}[H]
\centering
\caption{Trazabilidad entre objetivos específicos y capítulos de la memoria.}
\label{tab:trazabilidad}
\begin{tabular}{@{}lll@{}}
\toprule
\textbf{Objetivo} & \textbf{Descripción breve} & \textbf{Capítulo} \\
\midrule
O1 & Estado del arte            & \ref{cap:estado-arte} \\
O2 & Preparación de datos       & \ref{cap:metodologia}, \ref{cap:implementacion} \\
O3 & Descriptores topológicos   & \ref{cap:fundamentos}, \ref{cap:metodologia} \\
O4 & Clasificadores sobre TDA   & \ref{cap:metodologia}, \ref{cap:resultados} \\
O5 & Comparación con DL         & \ref{cap:resultados} \\
O6 & Visualización              & \ref{cap:implementacion}, \ref{cap:resultados} \\
\bottomrule
\end{tabular}
\end{table}
